\chapter{はじめに}
\label{cha:introduction}

\section{研究背景}
近年，コンピュータグラフィックス技術の発展に伴い，
人物の顔を対象とした写実的な画像生成が
映画，ゲーム，VR/AR などの分野で広く用いられている。
これらの応用において，
人間の皮膚の質感をどれだけ現実に近づけられるかは，
映像全体のリアリティを左右する重要な要素である。

皮膚の見えに大きな影響を与える要因の一つとして，
皮膚内部で生じる表面下散乱が挙げられる。
表面下散乱は，
光が皮膚表面から内部に侵入し，
内部で複数回散乱した後，
異なる位置から出射する現象であり，
皮膚特有の柔らかさや透明感を生み出す。
この現象を表現するため，
従来から双方向表面下散乱反射率分布関数
(BSSRDF) を用いた研究が行われてきた。

一方で，人間の顔には皮膚だけでなく，
眉毛や髭といった毛髪構造が局所的に存在している。
特に髭は，
口周りや顎といった限られた領域に高密度で分布し，
皮膚表面付近で光を吸収・遮蔽する。
そのため，髭が存在する領域では，
皮膚のみの場合とは異なる
光のにじみや明度分布が観測される。

\section{研究課題}
従来の皮膚表面下散乱モデルの多くは，
皮膚を一様な媒質として近似しており，
髭のような局所的な構造物の影響は
明示的には考慮されていない。
その結果，
髭の生えている領域においても
皮膚のみの場合と同様の散乱特性が適用され，
見えが不自然になる場合がある。

髭の影響を厳密に扱うためには，
毛髪一本単位のジオメトリや，
異方的な散乱特性を考慮した
高度なモデルを導入する必要がある。
しかし，そのような手法は
計算コストが高く，
実用的なレンダリングへの適用が難しい。

したがって，
髭の存在による影響を
過度な計算コストを増加させることなく，
皮膚表面下散乱モデルに取り込む手法が求められている。

\section{本研究の目的と貢献}
本研究の目的は，
皮膚表面下散乱モデルに
髭の分布情報を導入することで，
髭の存在による見えの変化を
より現実的に表現することである。

具体的には，
髭の分布を髭密度として定義し，
この密度に基づいて
皮膚モデルと髭モデルを混合する
表面下散乱モデルを提案する。
本手法では，
髭をジオメトリとして直接扱うことなく，
密度マップとして集約することで，
計算効率と表現力の両立を図る。

本研究の主な貢献は以下の通りである。
\begin{itemize}
  \item 髭密度に基づいて皮膚表面下散乱特性を制御する混合モデルの提案
  \item 髭分布をテクスチャとして表現し，
        レンダリング時に参照可能な密度マップ生成手法の提示
  \item 提案手法をモンテカルロパストレーシングに統合し，
        視覚的な効果と計算効率の両立を確認
\end{itemize}

\section{本論文の構成}
本論文は以下の構成からなる。
第2章では，
皮膚表面下散乱および関連する既存研究について述べる。
第3章では，
本研究で提案する
髭密度を考慮した皮膚表面下散乱モデルについて詳述する。
第4章では，
提案手法の評価方法および実験条件を示す。
第5章では，
実験結果を示し，結果に基づく考察を行う。
第6章で本研究のまとめと今後の課題を述べる。

